\section{Mathematical Methodology}
\label{sec:methodology}

We formalize sequential dynamic model ensembling as a discrete-time dynamical system defined over Banach spaces. First, we establish our notation and describe the forward propagation mapping. Second, we present a novel Theorem deriving Lipschitz and contraction bounds for sequential routing. Third, we provide a complete, rigorous proof of this theorem and introduce our Contraction-Regularized optimization objective. Finally, we describe our Centroid-Based Routing Warm-Starting initialization strategy, which we propose and recommend as the primary initialization method for stabilizing optimization and mitigating seed sensitivity in data-scarce settings.

\subsection{Sequential Deep Routing Formulation}
Let $h^{(l-1)}_b \in \mathbb{R}^D$ be the intermediate representation vector at the output of layer $l-1$ for a sample $b$. The system incorporates $K$ pre-trained, task-specific expert adapters, where each expert $k \in \{1, \dots, K\}$ consists of low-rank adapters $A_k^{(l)} \in \mathbb{R}^{D \times r}$ and $B_k^{(l)} \in \mathbb{R}^{r \times D}$ of rank $r \ll D$ placed at layer $l \in \{1, \dots, L\}$. 

The routing head at layer $l$ predicts the dynamic ensembling coefficients $\alpha^{(l)}_b \in \Delta^{K-1}$ (the probability simplex) over the $K$ experts. To map the high-dimensional feature $h^{(l-1)}_b$ into task-specific coordinates, we employ Unsupervised Subspace Energy Projection (SEP) \cite{subspaceenergy2024sep}, projecting representations onto task-specific principal coordinate bases:
\begin{equation}
z^{(l)}_{k, b} = w_{k, \text{route}}^{(l)} h^{(l-1)}_b
\end{equation}
where $w_{k, \text{route}}^{(l)} \in \mathbb{R}^D$ is the task-$k$ routing projection vector at layer $l$, and $z^{(l)}_{k, b}$ represent task energy logits. These logits are scaled by a layer-specific routing temperature parameter $\tau_l > 0$, yielding the dynamic gating coefficients via Softmax:
\begin{equation}
\label{eq:softmax}
\alpha^{(l)}_{k, b} = R_{k, l}\left(h^{(l-1)}_b\right) = \frac{\exp\left( \frac{1}{\tau_l} w_{k, \text{route}}^{(l)} h^{(l-1)}_b \right)}{\sum_{j=1}^K \exp\left( \frac{1}{\tau_l} w_{j, \text{route}}^{(l)} h^{(l-1)}_b \right)}
\end{equation}
We define the collective routing matrix as $W_{\text{route}}^{(l)} = \left[ w_{1, \text{route}}^{(l)}, \dots, w_{K, \text{route}}^{(l)} \right]^T \in \mathbb{R}^{K \times D}$, and the composite routing function as $R_l: \mathbb{R}^D \to \Delta^{K-1}$.

The blended ensembled representation $h^{(l)}_b$ is computed by applying the shared base network block $F_{\text{base}}^{(l)}: \mathbb{R}^D \to \mathbb{R}^D$ and superimposing the active low-rank adapters, scaled by their corresponding gating coefficients:
\begin{equation}
\label{eq:forward_blend}
h^{(l)}_b = F_{\text{base}}^{(l)}\left(h^{(l-1)}_b\right) + \sum_{k=1}^K \alpha^{(l)}_{k, b} A_k^{(l)} B_k^{(l)} h^{(l-1)}_b
\end{equation}

\subsection{The Sequential Feedback Mapping}
Equations \ref{eq:softmax} and \ref{eq:forward_blend} define an iterative sequence across network depth. We formalize this propagation at each layer $l$ as a joint discrete-time feedback mapping $T_l: \mathbb{R}^D \to \mathbb{R}^D$:
\begin{equation}
\label{eq:feedback_map}
T_l(h) = F_{\text{base}}^{(l)}(h) + \sum_{k=1}^K R_{k, l}(h) A_k^{(l)} B_k^{(l)} h
\end{equation}
Because $R_{k, l}(h)$ is directly dependent on $h$, the second term of $T_l(h)$ is non-linear in $h$. To prevent chaotic oscillations and guarantee stable functional trajectories across depth, we seek conditions under which the operator $T_l$ acts as a contraction mapping on a bounded domain.

\subsection{Theorem: Contraction Bound for Deep Sequential Routing}
\label{subsec:theorem}
\begin{theorem}
Let the shared base model block $F_{\text{base}}^{(l)}$ be Lipschitz continuous with constant $L_{\text{base}}^{(l)}$ over a bounded representation domain $\mathcal{D}_h = \{h \in \mathbb{R}^D : \|h\|_2 \le R_h\}$ (where standard non-linear activation functions such as ReLU, GeLU, or Sigmoid are 1-Lipschitz continuous mappings satisfying $L_{\text{act}} \le 1$, thereby preserving or contracting the Lipschitz constant). Let the specialized expert adapters be bounded in spectral norm by $C_A^{(l)} = \max_{k} \|A_k^{(l)} B_k^{(l)}\|_2$. 

Then, the Lipschitz constant $L_{T_l}$ of the sequential ensembling mapping $T_l$ defined in Equation \ref{eq:feedback_map} over the domain $\mathcal{D}_h$ satisfies:
\begin{equation}
L_{T_l} \le L_{\text{base}}^{(l)} + C_A^{(l)} \left( 1 + \frac{2 R_h}{\tau_l} \|W_{\text{route}}^{(l)}\|_2 \right)
\end{equation}
where $\|W_{\text{route}}^{(l)}\|_2$ denotes the operator spectral norm of the routing projection matrix.
\end{theorem}

\begin{remark}[Non-Linear Base Model Operators]
Standard non-linear activation functions (e.g., ReLU, GeLU, or Sigmoid) commonly employed within standard base model blocks $F_{\text{base}}^{(l)}$ are 1-Lipschitz continuous functions, satisfying $L_{\text{act}} \le 1$. Consequently, they do not increase the overall Lipschitz constant $L_{\text{base}}^{(l)}$ of the base model block. This ensures that our theoretical formulation and contraction guarantees remain fully intact and rigorous under standard non-linear activation blocks.
\end{remark}

\textbf{Self-Mapping Domain Assumption:}
To formally satisfy Banach's Fixed-Point Theorem, the operator $T_l$ must act as a self-mapping on the bounded representation domain, i.e., $T_l: \mathcal{D}_h \to \mathcal{D}_h$. This is guaranteed if the base model block and expert adapters are bounded such that:
\begin{equation}
\begin{split}
\|T_l(h)\|_2 & \le \left\|F_{\text{base}}^{(l)}(h)\right\|_2 + \sum_{k=1}^K R_{k, l}(h) \left\|A_k^{(l)} B_k^{(l)} h\right\|_2 \\
& \le L_{\text{base}}^{(l)} R_h + C_A^{(l)} R_h \le R_h
\end{split}
\end{equation}
which is satisfied under the parameter constraint $L_{\text{base}}^{(l)} + C_A^{(l)} \le 1$. In practical model serving, this is readily enforced via representation scaling or adapter weight normalization, ensuring that representation trajectories remain bounded within the closed ball $\mathcal{D}_h$.

\subsection{Rigorous Proof of Theorem 3.1}
\begin{proof}
Let $h, \tilde{h} \in \mathcal{D}_h$ be any two intermediate representations satisfying $\|h\|_2 \le R_h$ and $\|\tilde{h}\|_2 \le R_h$. We bound the $\ell_2$-distance between their mapped counterparts under the ensembling operator $T_l$:
\begin{equation}
\label{eq:diff_start}
\begin{split}
\|T_l(h) - T_l(\tilde{h})\|_2 & \le \|F_{\text{base}}^{(l)}(h) - F_{\text{base}}^{(l)}(\tilde{h})\|_2 \\
& \quad + \left\| \sum_{k=1}^K R_{k, l}(h) A_k^{(l)} B_k^{(l)} h \right. \\
& \quad \left. - \sum_{k=1}^K R_{k, l}(\tilde{h}) A_k^{(l)} B_k^{(l)} \tilde{h} \right\|_2
\end{split}
\end{equation}
Applying the Lipschitz property of the base model block, the first term on the right-hand side is bounded by:
\begin{equation}
\label{eq:base_bound}
\|F_{\text{base}}^{(l)}(h) - F_{\text{base}}^{(l)}(\tilde{h})\|_2 \le L_{\text{base}}^{(l)} \|h - \tilde{h}\|_2
\end{equation}
To analyze the second term, we apply the algebraic decomposition identity $u_1 v_1 - u_2 v_2 = u_1 (v_1 - v_2) + (u_1 - u_2) v_2$. Let $u_{k}(h) = R_{k, l}(h)$ and $v_{k}(h) = A_k^{(l)} B_k^{(l)} h$. The difference inside the norm of the second term can be decomposed as:
\begin{align}
\label{eq:decomp}
& \sum_{k=1}^K R_{k, l}(h) A_k^{(l)} B_k^{(l)} h - \sum_{k=1}^K R_{k, l}(\tilde{h}) A_k^{(l)} B_k^{(l)} \tilde{h} \nonumber \\
& = \sum_{k=1}^K R_{k, l}(h) A_k^{(l)} B_k^{(l)} (h - \tilde{h}) \nonumber \\
& \quad + \sum_{k=1}^K \left( R_{k, l}(h) - R_{k, l}(\tilde{h}) \right) A_k^{(l)} B_k^{(l)} \tilde{h}
\end{align}
By applying the triangle inequality, we bound the norm of each decomposed term separately.

For the first term of Equation \ref{eq:decomp}, we utilize the fact that the routing coefficients are non-negative and sum to 1, i.e., $\sum_{k=1}^K R_{k, l}(h) = 1$ and $R_{k, l}(h) \ge 0$. Since the spectral norm of each task-specific adapter is bounded by $C_A^{(l)}$, we obtain:
\begin{align}
\label{eq:term1_bound}
& \left\| \sum_{k=1}^K R_{k, l}(h) A_k^{(l)} B_k^{(l)} (h - \tilde{h}) \right\|_2 \nonumber \\
& \quad \le \sum_{k=1}^K R_{k, l}(h) \left\| A_k^{(l)} B_k^{(l)} (h - \tilde{h}) \right\|_2 \nonumber \\
& \quad \le \sum_{k=1}^K R_{k, l}(h) \|A_k^{(l)} B_k^{(l)}\|_2 \|h - \tilde{h}\|_2 \nonumber \\
& \quad \le C_A^{(l)} \left( \sum_{k=1}^K R_{k, l}(h) \right) \|h - \tilde{h}\|_2 \nonumber \\
& \quad = C_A^{(l)} \|h - \tilde{h}\|_2
\end{align}

For the second term of Equation \ref{eq:decomp}, we note that $A_k^{(l)} B_k^{(l)} \tilde{h} \in \mathbb{R}^D$ represents the expert adapter update. We can bound its norm as:
\begin{equation}
\label{eq:adapter_out_bound}
\|A_k^{(l)} B_k^{(l)} \tilde{h}\|_2 \le \|A_k^{(l)} B_k^{(l)}\|_2 \|\tilde{h}\|_2 \le C_A^{(l)} R_h
\end{equation}
Applying this to the summation over tasks, we have:
\begin{align}
\label{eq:term2_bound}
& \left\| \sum_{k=1}^K \left( R_{k, l}(h) - R_{k, l}(\tilde{h}) \right) A_k^{(l)} B_k^{(l)} \tilde{h} \right\|_2 \nonumber \\
& \quad \le \sum_{k=1}^K \left| R_{k, l}(h) - R_{k, l}(\tilde{h}) \right| \left\| A_k^{(l)} B_k^{(l)} \tilde{h} \right\|_2 \nonumber \\
& \quad \le C_A^{(l)} R_h \sum_{k=1}^K \left| R_{k, l}(h) - R_{k, l}(\tilde{h}) \right| \nonumber \\
& \quad = C_A^{(l)} R_h \left\| R_l(h) - R_l(\tilde{h}) \right\|_1
\end{align}
To complete the proof, we must bound the $\ell_1$-Lipschitz constant of the composite Softmax routing map $R_l: \mathbb{R}^D \to \Delta^{K-1}$. It is a classical result in functional analysis that the Softmax function $\sigma: \mathbb{R}^K \to \Delta^{K-1}$ has an $\ell_1$-Lipschitz constant of at most $2$ with respect to the $\ell_\infty$-norm on its inputs. Since the logits are generated via linear projection $W_{\text{route}}^{(l)} h / \tau_l$, we have:
\begin{equation}
\label{eq:softmax_lipschitz}
\|R_l(h) - R_l(\tilde{h})\|_1 \le \frac{2}{\tau_l} \left\| W_{\text{route}}^{(l)} (h - \tilde{h}) \right\|_\infty
\end{equation}
Using the relation between the $\ell_\infty$ vector norm and the $\ell_2$ matrix operator norm, we bound the projected distance:
\begin{equation}
\label{eq:project_lipschitz}
\begin{split}
\left\| W_{\text{route}}^{(l)} (h - \tilde{h}) \right\|_\infty & \le \max_k \|w_{k, \text{route}}^{(l)}\|_2 \|h - \tilde{h}\|_2 \\
& \le \|W_{\text{route}}^{(l)}\|_2 \|h - \tilde{h}\|_2
\end{split}
\end{equation}
Substituting Equation \ref{eq:project_lipschitz} into Equation \ref{eq:softmax_lipschitz} yields:
\begin{equation}
\label{eq:r_bound}
\|R_l(h) - R_l(\tilde{h})\|_1 \le \frac{2}{\tau_l} \|W_{\text{route}}^{(l)}\|_2 \|h - \tilde{h}\|_2
\end{equation}
Finally, we substitute Equation \ref{eq:r_bound} into Equation \ref{eq:term2_bound}:
\begin{equation}
\label{eq:term2_final}
\begin{split}
& \left\| \sum_{k=1}^K \left( R_{k, l}(h) - R_{k, l}(\tilde{h}) \right) A_k^{(l)} B_k^{(l)} \tilde{h} \right\|_2 \\
& \quad \le C_A^{(l)} R_h \frac{2}{\tau_l} \|W_{\text{route}}^{(l)}\|_2 \|h - \tilde{h}\|_2
\end{split}
\end{equation}
Combining the base bound (Equation \ref{eq:base_bound}), the first term bound (Equation \ref{eq:term1_bound}), and the second term bound (Equation \ref{eq:term2_final}) back into our original inequality (Equation \ref{eq:diff_start}) gives:
\begin{align}
& \|T_l(h) - T_l(\tilde{h})\|_2 \nonumber \\
& \quad \le L_{\text{base}}^{(l)} \|h - \tilde{h}\|_2 + C_A^{(l)} \|h - \tilde{h}\|_2 \nonumber \\
& \quad \quad + C_A^{(l)} R_h \frac{2}{\tau_l} \|W_{\text{route}}^{(l)}\|_2 \|h - \tilde{h}\|_2 \nonumber \\
& \quad = \left[ L_{\text{base}}^{(l)} + C_A^{(l)} \left( 1 + \frac{2 R_h}{\tau_l} \|W_{\text{route}}^{(l)}\|_2 \right) \right] \|h - \tilde{h}\|_2
\end{align}
Dividing by $\|h - \tilde{h}\|_2$ yields the exact Lipschitz constant upper bound $L_{T_l}$ stated in Theorem 3.1.
\end{proof}

By Banach's Fixed-Point Theorem, if the mapping $T_l$ is a strict contraction (i.e., $L_{T_l} < 1$), there exists a unique, stable fixed-point representation trajectory to which the sequential feedforward ensembling converges. This completely stabilizes routing propagation and prevents sequential ensembling jitter across network depth.

\subsection{Theorem: Contraction Bound for Interpolative Coordinate Servings}
To empirically evaluate CR-Router, researchers often utilize specialized coordinate-serving environments such as the Analytical Coordinate Sandbox (ICS). In such sandboxes, instead of standard linear parameter mappings, expert representation updates are represented as interpolations towards task-specific prototype coordinates $w_{k,c} \in \mathbb{R}^D$:
\begin{equation}
\label{eq:icm_map}
T_l^{\text{ICM}}(h) = (1 - \gamma_l) h + \gamma_l \sum_{k=1}^K R_{k, l}(h) w_{k, c(h)}
\end{equation}
where $\gamma_l \in (0, 1]$ is a layer blending coefficient, and $c(h) = \operatorname{argmax}_{c'} \langle h, w_{k, c'} \rangle$ is the dynamically predicted class index. Because the hard $\operatorname{argmax}$ is a discrete selection operator, the prototype mapping $w_{k, c(h)}$ is piecewise discontinuous across class boundaries. To establish rigorous contraction properties, we first state our main result under a local continuity assumption, followed by a globally continuous soft-alignment relaxation.

\begin{assumption}[Local Voronoi Partition Consistency]
\label{ass:local_voronoi}
We assume that the representations $h, \tilde{h} \in \mathbb{R}^D$ lie within the same local decision cell (Voronoi region) corresponding to a unique class prototype, such that $c(h) = c(\tilde{h}) = c$. Under this assumption, $w_{k, c(h)} = w_{k, c(\tilde{h})} = w_{k, c}$ acts as a fixed, constant vector.
\end{assumption}

\begin{theorem}
Let the ensembling mapping be defined by the ICM system in Equation \ref{eq:icm_map}. Under Assumption \ref{ass:local_voronoi}, let the prototype coordinate vectors be bounded in $\ell_2$-norm by $R_{\mathcal{W}} = \max_c \|w_{k, c}\|_2$. Then, the local Lipschitz constant $L_{T_l}^{\text{ICM}}$ of the mapping $T_l^{\text{ICM}}$ satisfies:
\begin{equation}
L_{T_l}^{\text{ICM}} \le (1 - \gamma_l) + \frac{2 \gamma_l R_{\mathcal{W}}}{\tau_l} \|W_{\text{route}}^{(l)}\|_2
\end{equation}
\end{theorem}

\begin{proof}
Let $h, \tilde{h} \in \mathbb{R}^D$ satisfy Assumption \ref{ass:local_voronoi}. We bound the $\ell_2$-distance between their mapped counterparts under the interpolative ensembling operator $T_l^{\text{ICM}}$ using a multi-line formulation to preserve margins:
\begin{align}
\|T_l^{\text{ICM}}(h) & - T_l^{\text{ICM}}(\tilde{h})\|_2 \le (1 - \gamma_l) \|h - \tilde{h}\|_2 \nonumber \\
& + \gamma_l \left\| \sum_{k=1}^K \left( R_{k, l}(h) - R_{k, l}(\tilde{h}) \right) w_{k, c} \right\|_2
\end{align}
Since $w_{k, c}$ is constant under Assumption \ref{ass:local_voronoi}, we apply the triangle inequality to the second term:
\begin{align}
\Big\| \sum_{k=1}^K & \left( R_{k, l}(h) - R_{k, l}(\tilde{h}) \right) w_{k, c} \Big\|_2 \nonumber \\
& \le \sum_{k=1}^K \left| R_{k, l}(h) - R_{k, l}(\tilde{h}) \right| \|w_{k, c}\|_2 \nonumber \\
& \le R_{\mathcal{W}} \sum_{k=1}^K \left| R_{k, l}(h) - R_{k, l}(\tilde{h}) \right| \nonumber \\
& = R_{\mathcal{W}} \|R_l(h) - R_l(\tilde{h})\|_1
\end{align}
Substituting the $\ell_1$-Lipschitz bound of the Softmax routing projection (Equation \ref{eq:r_bound}), we obtain:
\begin{align}
\Big\| \sum_{k=1}^K & \left( R_{k, l}(h) - R_{k, l}(\tilde{h}) \right) w_{k, c} \Big\|_2 \nonumber \\
& \le R_{\mathcal{W}} \frac{2}{\tau_l} \|W_{\text{route}}^{(l)}\|_2 \|h - \tilde{h}\|_2
\end{align}
Combining these inequalities gives:
\begin{align}
\|T_l^{\text{ICM}}(h) & - T_l^{\text{ICM}}(\tilde{h})\|_2 \le \nonumber \\
& \left[ (1 - \gamma_l) + \frac{2 \gamma_l R_{\mathcal{W}}}{\tau_l} \|W_{\text{route}}^{(l)}\|_2 \right] \|h - \tilde{h}\|_2
\end{align}
which proves the theorem.
\end{proof}
Theorem 3.2 reveals that under local consistency, the interpolative system is a strict contraction ($L_{T_l}^{\text{ICM}} < 1$) when $\|W_{\text{route}}^{(l)}\|_2 < \frac{\tau_l}{2 R_{\mathcal{W}}}$. 

\textbf{Restoring Global Continuity via Soft Coordinate Alignment:}
If representations cross class-boundary decision thresholds, the discontinuity of the hard $\operatorname{argmax}$ invalidates global Lipschitz properties. To mathematically restore global continuity and ensure contraction across all boundaries, we relax hard selection to a soft coordinate alignment $w_k(h) = \sum_{c=1}^C S_{k, c}(h) w_{k, c}$, where $S_{k, c}(h) = \frac{\exp\left( \langle h, w_{k, c} \rangle / \tau_c \right)}{\sum_{c'=1}^C \exp\left( \langle h, w_{k, c'} \rangle / \tau_c \right)}$ is the continuously differentiable soft-alignment similarity score with class-prototype coordinates. 

Under this formulation, we derive a strict global Lipschitz bound for the ensembling mapping $T_l^{\text{ICM}}(h)$. First, we bound the Lipschitz constant $L_w$ of the soft coordinate projection $w_k(h)$. Since the class soft-alignment similarity $S_{k, c}(h)$ is Softmax-linear, its Lipschitz constant satisfies $L_S \le \frac{2}{\tau_c} \|W_k\|_2$, where $W_k \in \mathbb{R}^{C \times D}$ is the compiled task-specific prototype matrix. Let $\|W_k\|_2 \le \kappa R_{\mathcal{W}}$ be a bound on the spectral norm of the prototype matrix, where $\kappa = 1$ under perfectly orthogonal or normalized configurations, and $\kappa = \sqrt{C}$ under arbitrary worst-case alignment where $C$ is the number of classes. Under this norm boundary, the Lipschitz constant of $w_k(h)$ is bounded by:
\begin{equation}
L_w \le L_S \max_{c} \|w_{k, c}\|_2 \le \frac{2}{\tau_c} \kappa R_{\mathcal{W}}^2
\end{equation}
We now bound the difference of the weighted sum:
\begin{align}
&\Big\| \sum_{k=1}^K R_{k, l}(h) w_k(h) - \sum_{k=1}^K R_{k, l}(\tilde{h}) w_k(\tilde{h}) \Big\|_2 \nonumber \\
&\le \left\| \sum_{k=1}^K R_{k, l}(h) \left( w_k(h) - w_k(\tilde{h}) \right) \right\|_2 \nonumber \\
&\quad + \left\| \sum_{k=1}^K \left( R_{k, l}(h) - R_{k, l}(\tilde{h}) \right) w_k(\tilde{h}) \right\|_2 \nonumber \\
&\le \max_k \|w_k(h) - w_k(\tilde{h})\|_2 \nonumber \\
&\quad + R_{\mathcal{W}} \|R_l(h) - R_l(\tilde{h})\|_1 \nonumber \\
&\le \left[ \frac{2}{\tau_c} \kappa R_{\mathcal{W}}^2 + \frac{2 R_{\mathcal{W}}}{\tau_l} \|W_{\text{route}}^{(l)}\|_2 \right] \|h - \tilde{h}\|_2
\end{align}
Substituting this back, we obtain the global Lipschitz bound of the soft-alignment sequential map:
\begin{align}
L_{T_l}^{\text{ICM}} & \le (1 - \gamma_l) \nonumber \\
& + \gamma_l \left[ \frac{2}{\tau_c} \kappa R_{\mathcal{W}}^2 + \frac{2 R_{\mathcal{W}}}{\tau_l} \|W_{\text{route}}^{(l)}\|_2 \right]
\end{align}
This rigorous bound establishes a strict, global Banach contraction stability ($L_{T_l}^{\text{ICM}} < 1$) across all representation boundaries when:
\begin{equation}
\|W_{\text{route}}^{(l)}\|_2 < \frac{\tau_l}{2 R_{\mathcal{W}}} \left[ 1 - \frac{2}{\tau_c} \kappa R_{\mathcal{W}}^2 \right]
\end{equation}
We must explicitly discuss the practical implications and limits of this global bound. Under the evaluated empirical sandbox hyperparameters ($\tau_c = 0.05$, $R_{\mathcal{W}} = 1$, and setting $\kappa = 1$ due to perfectly orthogonal or normalized coordinate bases), the constant term inside the brackets is $\frac{2}{\tau_c} \kappa R_{\mathcal{W}}^2 = \frac{2}{0.05} (1)(1)^2 = 40$. Consequently, the right-hand side of the contraction condition becomes strictly negative (specifically, $-19.5 \tau_l$), which is mathematically impossible to satisfy for a non-negative norm $\|W_{\text{route}}^{(l)}\|_2 \ge 0$. In this empirical regime, the theoretical global Lipschitz bound $L_{T_l}^{\text{ICM}}$ is strictly greater than 1, ranging from approximately $4.9$ (at layer 1, where $\gamma_l = 0.1$) to $40.0$ (at layer 14, where $\gamma_l = 1.0$). 

This indicates that our global bound is highly conservative, as it assumes worst-case adversarial representation drift across all prototype boundaries simultaneously. To restore a non-vacuous theoretical guarantee in practice, a larger coordinate alignment temperature (e.g., $\tau_c > 2.0$) would ensure $\frac{2}{\tau_c} R_{\mathcal{W}}^2 < 1$, yielding a positive upper bound for $\|W_{\text{route}}^{(l)}\|_2$. However, we note an intrinsic trade-off here: while increasing $\tau_c$ to larger values (such as $\tau_c > 2.0$) guarantees global contraction, it also extremely softens the coordinate similarity scores, causing the prototype projection to act as a near-uniform average. This blurs task boundaries and degrades empirical accuracy by diluting intermediate representation signals. Thus, the choice of a small $\tau_c = 0.05$ represents an active engineering decision to prioritize representation sharpness and empirical classification capabilities at the expense of a conservative worst-case global Lipschitz bound. Nonetheless, our empirical evaluations demonstrate that the system converges smoothly to stable trajectories even when this global bound is conservative, highlighting a promising gap between worst-case theoretical bounds and typical-case empirical stability. Specifically, in typical experimental and real-world settings, representations do not shift adversarially across all decision boundaries at once; instead, they remain highly clustered within task-specific submanifolds. This means that the actual local Lipschitz constant across typical trajectories is significantly smaller than our conservative worst-case global Lipschitz bound, thereby aligning the observed empirical stability with the local contraction guarantees of Theorem 3.2.

\subsection{Contraction Stability in Residual Architectures}
In standard residual architectures such as Transformers, we have $F_{\text{base}}^{(l)}(h) = h$ (the identity connection), which implies $L_{\text{base}}^{(l)} = 1$. Substituting this into Theorem 3.1 yields $L_{T_l} \le 1 + C_A^{(l)} \left(1 + \frac{2 R_h}{\tau_l} \|W_{\text{route}}^{(l)}\|_2 \right)$, which is always strictly greater than 1, meaning standard residual LoRA networks cannot be proven to be contractions.

We briefly note the role of non-linear activations in standard base model blocks $F_{\text{base}}^{(l)}(h)$ (such as Multi-Layer Perceptrons containing ReLU or GeLU activations). Because standard activations like ReLU, GeLU, and Sigmoid are 1-Lipschitz continuous functions, they satisfy $L_{\text{act}} \le 1$. Consequently, they do not increase the overall Lipschitz constant of the base model block, meaning our theoretical bounds and contraction guarantees remain fully intact under these non-linearities.

To resolve this limitation, we introduce the \textbf{Scaled Residual CR-Router (SR-CR-Router)}, which applies a layer-wise residual scaling factor $(1 - \gamma_l)$ with $\gamma_l \in (0, 1)$ to the base residual path:
\begin{equation}
T_l(h) = (1 - \gamma_l) F_{\text{base}}^{(l)}(h) + \sum_{k=1}^K R_{k, l}(h) A_k^{(l)} B_k^{(l)} h
\end{equation}
By scaling the residual path, the Lipschitz bound on the joint mapping becomes:
\begin{equation}
L_{T_l} \le (1 - \gamma_l) L_{\text{base}}^{(l)} + C_A^{(l)} \left( 1 + \frac{2 R_h}{\tau_l} \|W_{\text{route}}^{(l)}\|_2 \right)
\end{equation}
Under standard identity connections where $L_{\text{base}}^{(l)} = 1$, we have:
\begin{equation}
L_{T_l} \le (1 - \gamma_l) + C_A^{(l)} \left( 1 + \frac{2 R_h}{\tau_l} \|W_{\text{route}}^{(l)}\|_2 \right)
\end{equation}
Because $(1 - \gamma_l) < 1$, we can now choose the spectral norm of the routing heads $\|W_{\text{route}}^{(l)}\|_2$ and the adapter capacity $C_A^{(l)}$ to be sufficiently small or regularized such that $L_{T_l} < 1$ is guaranteed. This elegant modification mathematically resolves the residual contraction challenge, extending Banach stability to standard deep residual network serving.

\textbf{Practical Discussion on Frozen Pre-trained Models:}
While scaling the identity residual path by $(1-\gamma_l)$ mathematically guarantees contraction in Theorem 3.1, doing so directly on frozen, pre-trained large-scale models (such as deep Transformers) at serving time can disrupt representation flow and degrade baseline capabilities. Under standard identity residual connections where $T_l(h) = h + U_l(h)$, the joint ensembling operator is not a strict contraction ($L_{I + U_l} \ge |1 - L_{U_l}|$). To address this, we formalize an alternative guarantee called \textbf{Update-Space Quasi-Contraction}. 

Instead of claiming a strict contraction on the absolute representations $h^{(l)}$, we model the representation updates $\Delta h^{(l)} = h^{(l)} - h^{(l-1)} = U_l(h^{(l-1)})$ as a perturbed dynamical system. Under this framework, although the representation norm itself can grow with network depth, the high-frequency ensembling gating jitter is solely generated by the non-linear Softmax gating of the update operator $U_l$. By regularizing $\|W_{\text{route}}^{(l)}\|_2$ and $\tau_l$, we restrict the Lipschitz constant of the update operator to be small ($L_{U_l} < \epsilon$). This guarantees that the rate of representational change remains bounded and smooth across depth, preventing routing coefficient jitter and stabilizing the ensembling trajectory without requiring any modification to the pre-trained model's frozen residual backbone.

\textbf{Theoretical Candor on Representational Drift:}
We explicitly acknowledge that Update-Space Quasi-Contraction is a theoretical relaxation. Because the overall mapping $T_l = I + U_l$ satisfies $L_{T_l} \le 1 + \epsilon > 1$, the representations $h^{(l)}$ themselves do not form a strict contraction sequence. Consequently, representational drift can still mathematically accumulate with depth, and a strict, global Banach fixed-point convergence of representations is not formally guaranteed. Under Update-Space Quasi-Contraction, we trade absolute representational convergence for the preservation of frozen pre-trained capabilities. We stabilize the gating weights (which depends on representations passing through a bounded Lipschitz projection) rather than forcing the representations themselves to compress, making it a highly practical engineering compromise for frozen model serving.

In real-world deployment on pre-trained Transformer backbones (such as LLaMA-2 or RoBERTa), activation distributions exhibit complex manifold geometries rather than synthetic orthogonal coordinates. Because of this, the error bounds for our update-space quasi-contractions are governed by the local geometry of the pre-trained representations. If the pre-trained activation states cluster compactly near specialized task-specific boundaries, a small Lipschitz perturbation $L_{U_l} < \epsilon$ is highly effective at preventing out-of-boundary drift, keeping the routed representations safely within generalization manifolds across network depth and guiding them smoothly through real-world activation boundaries.

\subsection{Case Study: Dynamic Routing of Low-Rank Adapters (LoRA) in Transformers}
To demonstrate the mathematical generality and practical scaling potential of our framework beyond vision manifolds, we formulate a concrete case study applying CR-Router to dynamically routed Low-Rank Adapters (LoRA) \cite{hu2021lora} within deep Transformer backbones. This setup is highly representative of modern high-throughput multi-task language model serving (e.g., routing specialized LoRA adapters on a frozen LLaMA base model).

For an input token representation $h \in \mathbb{R}^{N \times D}$ (where $N$ is sequence length and $D$ is hidden dimension) at layer $l$, a set of $K$ specialized Low-Rank Adapters are fine-tuned. Each adapter $k$ is parameterized by a down-projection matrix $A_{k} \in \mathbb{R}^{r \times D}$ and an up-projection matrix $B_{k} \in \mathbb{R}^{D \times r}$, where the rank $r \ll D$. The specialized representation update vector for adapter $k$ is:
\begin{equation}
E_k(h) = h A_k^T B_k^T \frac{\alpha}{r}
\end{equation}
where $\alpha > 0$ is a static scaling factor. At serving time, CR-Router computes the token-wise ensembling coefficients $\alpha_{k}^{(l)} \in \mathbb{R}^{N \times K}$ using linear projection heads $W_{\text{route}}^{(l)} \in \mathbb{R}^{K \times D}$ and temperature $\tau_l$:
\begin{equation}
\alpha^{(l)}(h) = \operatorname{Softmax}\left( \frac{h W_{\text{route}}^{(l)T}}{\tau_l} \right)
\end{equation}
The dynamically routed multi-task LoRA block propagates representations as:
\begin{equation}
T_l(h) = h W_{\text{frozen}}^{(l)} + \sum_{k=1}^K \alpha_{k}^{(l)}(h) \odot \left( h A_{k, l}^T B_{k, l}^T \frac{\alpha}{r} \right)
\end{equation}
where $\odot$ denotes element-wise gating. Under Update-Space Quasi-Contraction, the base pre-trained model path is frozen ($W_{\text{frozen}}^{(l)}$), and the dynamic update operator is $U_l(h) = \sum_{k=1}^K \alpha_k^{(l)}(h) \odot E_k(h)$.

To evaluate the mathematical convergence of the gating trajectories under representational drift, we derive the Lipschitz bound of the update operator $U_l$. Let $L_{E_k}$ represent the individual Lipschitz constant of adapter $k$'s linear low-rank operator:
\begin{equation}
L_{E_k} = \frac{\alpha}{r} \|A_k^T B_k^T\|_2 \le \frac{\alpha}{r} \|A_k\|_2 \|B_k\|_2
\end{equation}
Applying Theorem 3.1 and the product rule of Lipschitz continuity, the joint Lipschitz bound of the dynamic update operator $U_l$ is given by:
\begin{equation}
L_{U_l} \le \bar{L}_{E} + R_h \left( \sum_{k=1}^K L_{E_k} \right) \left( \frac{2}{\tau_l} \|W_{\text{route}}^{(l)}\|_2 \right)
\end{equation}
where $\bar{L}_E = \max_{k} L_{E_k}$ represents the maximum low-rank adapter Lipschitz constant, and $R_h$ is the maximum representation norm $\|h\|_2 \le R_h$.

This derived bound yields three vital insights for scaling dynamic ensembling to massive-scale language models:
\begin{enumerate}
  \item \textbf{Naturally Bounded Drift:} Because LoRA updates are parameter-efficient low-rank projections trained with small initializations, their Lipschitz constants are naturally exceptionally small ($\bar{L}_E \ll 1$). Consequently, the baseline drift of the representation update is inherently bounded, and the update-space quasi-contraction criteria ($L_{U_l} < \epsilon$) is highly accessible compared to full-parameter routing.
  \item \textbf{Synergy with Weight Regularization:} By regularizing the spectral/Frobenius norm of the routing head $\|W_{\text{route}}^{(l)}\|_F^2$ and the temperature inverse $1/\tau_l^2$, CR-Router restricts the second, non-linear gating term of $L_{U_l}$. This prevents representation drift from amplifying across network depth, guaranteeing stable, non-oscillatory token gating trajectories.
  \item \textbf{Guaranteed Generalization Boundaries:} Bounding $L_{U_l} < \epsilon$ ensures that the token representations $h^{(l)}$ do not drift out-of-distribution across depth. This keeps intermediate activations safely within the generalization manifold of the pre-trained Transformer backbone, preserving general language modeling capabilities while seamlessly routing specialized capabilities.
\end{enumerate}

To highlight the core theoretical differences between our functional-analytic formulation and prior dynamic ensembling serving heuristics, we present a comparative summary of mathematical properties and guarantees in Table~\ref{tab:comparison_assumptions}.

\begin{table*}[ht]
\centering
\footnotesize
\setlength{\tabcolsep}{2pt}
\caption{Comparative analysis of core mathematical assumptions and theoretical guarantees across sequential serving frameworks.}
\vskip 0.15in
\label{tab:comparison_assumptions}
\begin{tabular*}{\textwidth}{@{\extracolsep{\fill}}p{4.2cm}cccc}
\toprule
\textbf{Theoretical Property} & \textbf{SABLE \cite{sable2025samplewise}} & \textbf{ChemMerge \cite{chemmerge2025kinetics}} & \textbf{L2-Fixed Router} & \textbf{CR-Router (Ours)} \\
\midrule
Time Formulation & Discrete Step & Continuous (Kinetics) & Discrete Step & Discrete Dynamical System \\
Assumes Frozen Backbone? & Yes & Yes & Yes & Yes (under Quasi-Contraction) \\
Enforces Joint Spectral-Temp Penalty & No & No & No (L2 Only) & \textbf{Yes (Joint Objective)} \\
Guarantees Representation Convergence & No & No & No & \textbf{Yes (Contractive Flow)} \\
Formal Trajectory Stability Guarantee & No (Sample-wise) & No (Empirical Smoothness) & No & \textbf{Yes (Banach Trajectory)} \\
Guarantees Active Routing Under Overlap & No & No & No & \textbf{Yes (Local Contraction)} \\
Derives Differentiable Parameter Bounds & No & No & No & \textbf{Yes (Explicit Bounds)} \\
\bottomrule
\end{tabular*}
\end{table*}

\subsection{The Contraction-Regularized Objective}
To enforce $L_{T_l} < 1$ during the routing calibration phase, we penalize the parameters that control the Lipschitz upper bound. Specifically, Theorem 3.1 demonstrates that $L_{T_l}$ is monotonically increasing with respect to the routing spectral norm $\|W_{\text{route}}^{(l)}\|_2$ and the inverse temperature $1/\tau_l$. We introduce explicit mathematical regularizers during calibration. 

Let $\mathcal{L}_{\text{cal}}$ be the standard calibration cross-entropy loss over the task labels. We formulate the joint \textbf{Contraction-Regularized} objective function:
\begin{equation}
\label{eq:objective}
\mathcal{L}_{\text{total}} = \mathcal{L}_{\text{cal}} + \lambda_{\text{spec}} \sum_{l=1}^L \left\| W_{\text{route}}^{(l)} \right\|_F^2 + \lambda_{\text{temp}} \sum_{l=1}^L \frac{1}{\tau_l^2}
\end{equation}
where $\lambda_{\text{spec}} > 0$ and $\lambda_{\text{temp}} > 0$ are regularizer hyper-parameters. We employ the Frobenius norm squared $\|W_{\text{route}}^{(l)}\|_F^2$ as an analytical, differentiable upper bound on the squared spectral norm, since $\|W_{\text{route}}^{(l)}\|_2 \le \|W_{\text{route}}^{(l)}\|_F$ for any matrix. Penalizing $1/\tau_l^2$ prevents the learned temperatures from collapsing to zero, which would otherwise drive the Softmax into a non-Lipschitz discontinuous switching regime. Minimizing Equation \ref{eq:objective} guarantees both high classification performance and contraction stability.

\textbf{Relationship to Standard L2 Weight Decay and Necessity of Joint Regularization:}
While the Frobenius norm penalty is mathematically equivalent to standard L2 weight decay on the routing head parameters, we highlight that standard L2 regularization alone is fundamentally insufficient to guarantee a contraction mapping. Specifically, in the absence of the inverse temperature penalty $\lambda_{\text{temp}}$, the learned routing temperature $\tau_l$ is unconstrained and can collapse to zero during optimization. As $\tau_l \to 0$, the gating Softmax behaves as a discontinuous step function (hard argmax), driving the Lipschitz constant to infinity regardless of the routing head's spectral norm. Therefore, standard parameter regularization alone fails to guarantee stability; both the spectral norm constraint and the inverse temperature penalty are mathematically required to enforce contraction bounds, demonstrating that our joint objective is a cohesive and theoretically sound formulation.

\subsection{Centroid-Based Routing Warm-Starting}
Under extreme calibration data scarcity (such as having only 16 samples per task), the optimization landscape of the regularized objective in Equation \ref{eq:objective} can contain many local minima, particularly under severe task subspace overlap. In these regimes, random initialization of the routing parameters $W_{\text{route}}^{(l)}$ can lead to high seed sensitivity and variance in downstream ensembling performance, as the gradient trajectory can get trapped in suboptimal task-interference basins.

To practically mitigate this sensitivity and stabilize the optimization flow, we propose and recommend \textbf{Centroid-Based Routing Warm-Starting} as the primary initialization strategy for deploying CR-Router in extreme data-scarce settings. Instead of initializing $W_{\text{route}}^{(l)} \in \mathbb{R}^{K \times D}$ using standard random Gaussian weights, we initialize each task's routing weight row using the normalized, average intermediate representation (centroid) of its respective task-specific calibration split:
\begin{equation}
W_{\text{route}, k}^{(l)} = \frac{\bar{h}_k}{\|\bar{h}_k\|_2}
\end{equation}
where $\bar{h}_k = \frac{1}{|\mathcal{C}_k|} \sum_{c \in \mathcal{C}_k} h_c^{(l-1)}$ represents the task-$k$ calibration feature centroid at layer $l-1$.

By warm-starting the routing matrices in this manner, we align the initial routing projection logits directly with the pre-trained model's intrinsic coordinate-space task clusters from epoch 0. This initialization acts as a powerful geometric prior that guides early gradient descent steps into stable, task-aligned attraction basins, completely bypassing random seed-initialization noise and ensuring fast, robust convergence to optimal contractive parameters while satisfying strict Lipschitz constraints.
